This section states how the project will be managed. Its content aligns with the project planning and project assessment/control activities of ISO/IEC/IEEE 16326:2019 \cite{ISO16326}, set within the life-cycle process frameworks of ISO/IEC/IEEE 12207:2017 \cite{ISO12207} and 15288:2023 \cite{ISO15288}. Keep it concise; detailed recurring mechanics belong in the Documentation \& Reporting section.

\subsection{Life Cycle Model}
State the life cycle model the team will follow. Senior Design teams use an iterative, agile (Scrum) model \cite{Rubin2012, SchwaberSutherland2020}: fixed-length sprints, each producing a potentially shippable increment, spanning CSE 4316 (requirements and design) and CSE 4317 (implementation, test, and demonstration). Note the sprint length and the approximate number of sprints across both semesters (consult the course schedules).

\subsection{Project Processes \& Ceremonies}
Summarize the recurring processes the team will run: sprint planning, daily/weekly stand-ups, sprint review with the customer, and sprint retrospective. Identify who fills the Product Owner and Scrum Master roles and whether those roles rotate.

\subsection{Estimation \& Measurement}
Describe how work is estimated and how progress is measured and controlled --- the project assessment and control activity of ISO/IEC/IEEE 16326. State the unit of estimation (e.g. task-hours or story points), how the team tracks effort, and the primary progress metrics (sprint burndown, velocity, and requirement/test coverage). Explain the threshold at which a variance (e.g. burndown trending above the ideal line) triggers corrective action.

\subsection{Change \& Configuration Management}
State how requirement and scope changes are proposed, evaluated, and approved (who signs off), and how work products are version-controlled and baselined (e.g. Git repository, branching model, tagged document releases).
